In this section, we study the time and space complexity of computing the rank of a zero-one matrix monoid and finding a matrix of minimum rank in it. We provide two approaches with the same time complexity. The first one relies on the linear-algebraic tools developed in \autoref{sec-toolbox}. It turns out to have smaller space complexity than the second, combinatorial, approach, but is limited to only computing the rank. This is due to the fact that we never explicitly construct a maximal column in this approach, which is required in order to find a matrix of minimum rank by \autoref{thm:Cesari}. Moreover, it is not known if one can find a maximal column in \NC. In contrast, the combinatorial approach explicitly constructs a matrix of minimum rank step by step. In a way, this is exactly the reason why it requires more space to achieve the same time complexity: in the proof of \autoref{lemma:time-max-col}, we have to precompute a linear number of matrices (one for each step), since we do not know in advance which matrices we will need and since computing them ``on the fly'' would have a higher time complexity.
In the total DFAs case, where transition matrices require only linear amount of space and can be multiplied in linear time, the combinatorial approach becomes much more efficient, and outmatches its linear-algebraic counterpart by a factor of the alphabet size.

The three main results of this section are as follows.

\begin{thm}\label{thm:linalg-time}
    The (real) rank of an $n$-state UFA over an alphabet of size $m$ can be computed in $\OO(mn^4)$ time and $\OO(n^2)$ space. 
\end{thm}

\begin{thm}\label{thm:finding-matrix}
 A matrix of minimum (real) rank in an $n$-state UFA over an alphabet of size~$m$ can be found in  $\OO(mn^4)$ time and $\OO(n^3)$ space. 
\end{thm}

\begin{thm}\label{thm:main-dfas}
    A matrix of minimum (real) rank in an $n$-state total DFA over an alphabet of size~$m$ can be found in $\OO(n^3 + mn^2)$ time and $\OO(n^2)$ space.
\end{thm}

% Fix indentation
\noindent
Until the end of the section, fix a strongly connected complete UFA $\Aa = (Q, \Sigma, \Delta)$. Denote $n = |Q|$, $m = |\Sigma|$. \autoref{subsec:non-sc} shows that strong connectivity can be assumed without loss of generality.

\subsection{Square automaton and square digraph}\label{subsec:square-aut}

We will need the construction of the square automaton of an NFA. The \emph{square automaton} 
$\Aa^{(2)} = (Q^{(2)}, \Sigma, \Delta^{(2)})$ of $\Aa$ is defined as follows. Let $Q^{(2)} = \{(p, q) \mid p, q \in Q\}$, and for $p, q \in Q$ and $a \in \Sigma$, the transitions are defined component-wise, that is, 
\[\Delta^{(2)} = \{ ((p, q), a, (p', q')) \in Q^{(2)} \times \Sigma \times Q^{(2)} \mid (p, a, p'), (q,a,q') \in \Delta\}.\]
Note that the square automaton of a total DFA is also a total DFA. 


We call states of the form $(q, q)$ in $\Aa^{(2)}$ \emph{singletons}. Observe that the restriction of~$\Aa^{(2)}$ to singletons is equal to $\Aa$.
We denote by $G^{(2)} = (V^{(2)}, E^{(2)})$ the underlying digraph of~$\Aa^{(2)}$ obtained by forgetting the labels of the transitions. Note that $|E^{(2)}| = \OO(mn^4)$, and there exists an infinite series of complete UFAs over a two-letter alphabet with $|E^{(2)}| = \Theta(n^4)$~\cite[Appendix A]{Kiefer2019Arxiv}. If $\Aa$ is a total DFA, then $|E^{(2)}| = mn^2$.


\subsection{Minimum rank in \texorpdfstring{$\OO(mn^4)$}{O(n4)} time}\label{subsec:linalg-time}

We now perform the steps of the linear-algebraic algorithm described in \autoref{sec-nc2}, but implement them efficiently in terms of time and space complexity.

\begin{lem}\label{lem:mer-time}
    For a state $p$, the set $\Mer(p)$ can be computed in $\OO(mn^4)$ time and $\OO(n^2)$ space.
\end{lem}
\begin{proof}
    Perform a multi-source backwards digraph search starting from all singletons in $G^{(2)}$ and label all states $q \in Q$ such that $(p, q)$ or $(q, p)$ is visited during this search. This search can be performed in time linear in the number of edges of $|E^{(2)}|$, and $|E^{(2)}| = \OO(mn^4)$. 
    Observe that only the vertices of this digraph have to be stored in the memory explicitly, since the edges can be computed on the fly from the input without increasing the time complexity, hence the space complexity of the algorithm is $\OO(n^2)$.
\end{proof}

\begin{lem}\label{lem:pseudcolumn-time}
    A maximal pseudo-column can be found in $\OO(mn^4)$ time and $\OO(n^2)$ space.
\end{lem}
\begin{proof}
To use \autoref{thm:pseud-linsyst} we first need to set up the linear system of equations described there.
The average matrix $\overline{A}$ can be computed in time $\OO(m n^2)$.
The weight vectors $\alpha, \beta \in \mathbb{Q}^Q$ can then be computed in time $\OO(n^3)$ by solving a system of linear equations.
Then, using \autoref{lem:mer-time},
we compute in $\OO(mn^4)$ time the set $\Mer(p)$ for some state $p$.
We also need to compute a basis of the vector space~$U$ defined in \autoref{subsec:pseudocolumns}.
As in the proof of \autoref{lem:U-in-NC}, we compute a basis of $U = U' = \spn{\alpha^T M'(w) \mid w \in \Sigma^+}$, which is the smallest vector space that contains $\alpha^T M(a)$ for all $a \in \Sigma$ and is closed under post-multiplication with $M(a)$ for all $a \in \Sigma$.
This can be done in $\OO(mn^3)$ time and $\OO(n^2)$ space using a worklist algorithm and keeping a basis in echelon form using Gaussian elimination, as described, e.g., in \cite[Section 2]{KieferWeighted20}.
Finally, we solve the system of linear equations from \autoref{thm:pseud-linsyst}, 
which can be done in $\OO(n^3)$ time.
Each step requires at most $\OO(n^2)$ space.
\end{proof}

% Fix indentation
\noindent
By \autoref{lem:Cas}, we thus get that $\mcw$ and $\mrw$ can be computed in  $\OO(mn^4)$ time and $\OO(n^2)$ space, which together with \autoref{thm:rank-weight} proves \autoref{thm:linalg-time}. 
We remark that for total DFAs the proof of \autoref{lem:pseudcolumn-time} gives $\OO(mn^3)$ time for computing the rank. The combinatorial algorithm provided below will improve it to $\OO(n^3 + mn^2)$ (see \autoref{subsec-dfas}), while additionally finding a matrix of minimum rank.


\subsection{Efficiently constructing a maximal column}\label{subsec:max-col} 

We now consider a more general problem of finding a matrix of minimum rank in a zero-one matrix monoid. As mentioned above, by \autoref{thm:Cesari} we have to explicitly construct a maximal column for that, which turns out to be a more difficult task in terms of the time complexity. We will see in the next subsection that we only need one maximal column (together with a word constructing it) to get a matrix of minimum rank.
The goal of this subsection is thus to show how to efficiently compute a short representation of a word constructing a maximal column, since the word itself may be longer than the time complexity we are aiming for. We do so by reusing repeating subwords in such a word, and describing their occurrences with a straight line program.

In \cite[Section 5]{Eppstein1990}, an algorithm for constructing a word of rank one in total DFAs in $\OO(m^3 + mn^2)$ time and $\OO(n^2)$ space was suggested. Our approach for finding a maximal column and a word constructing it follows a similar direction, with a few key differences. Firstly, we observe that it is enough to only compute words merging states with one chosen state~$p$, instead of finding all pairs of mergeable states. This both simplifies the algorithm (in~\cite{Eppstein1990} an additional step is required to decrease the space complexity from $\OO(n^3)$ to $\OO(n^2)$, which we get for free) and allows to present our results in terms of straight line programs, giving a better insight into the regularities present in the constructed word of minimum rank.

We define set straight line programs (set-SLPs), which are just SLPs with multiple initial symbols, and thus encode a set of words instead of one word.
Formally, a \emph{set-SLP} is a tuple $(\VV, \Sigma, R, \Start)$, where $\VV$ and $\Sigma$ are disjoint finite sets of \emph{nonterminal} and \emph{terminal} symbols respectively, $R \colon \VV \to (\VV \cup \Sigma)^*$ is a function defining a derivation rule for each nonterminal symbol, and $\Start \subseteq \VV$ is a set of \emph{initial} symbols. For $v \in \VV$, we write $R(v)$ as $v \to w$ with $w \in (\VV \cup \Sigma)^*$, and we call $v$ and $w$ the left- and right-hand sides of this derivation rule respectively. The \emph{length} of a set-SLP is the total length of the right-hand sides of all the derivation rules.
The semantics of a set-SLP is defined as follows. Given an initial symbol $s \in \Start$, we recursively replace each symbol in $R(s)$ with the right-hand side of its derivation rule until we obtain a word over $\Sigma$, which is called \emph{the word encoded by~$s$}.
We require that each initial symbol produces a unique word over $\Sigma$ as a result of such derivation. Namely, we require that there exists a total linear order $\le$ on the set $\VV$ such that for all $v$ with $v \to w$, $w$ does not contain $v' \in \VV$ with $v' \le v$. The (multi-)set of words encoded by all initial symbols is called \emph{the set of words encoded by} a set-SLP.

\begin{exa}\label{ex:set-slp}
    Consider a set-SLP $(\{w_1, w_3, w_5, u_1, u_2, u_3\}, \{a, b\}, R, \{w_1, w_3, w_5\})$ with 
\[w_1 \to u_1, w_3 \to u_3 u_2, w_5 \to u_2, 
u_1 \to aab, u_2 \to aab, u_3 \to aaba.\]
This set-SLP encodes the (multi-)set $\{aab, aabaaab, aab\}$, and illustrates the reason why we are using set-SLPs: they allow to construct sets of words out of smaller ``pieces'' without having to explicitly repeat these ``pieces'' multiple times (in our example, we are reusing~$u_2$). Note that the set-SLPs that we construct below only encode sets of words whose total length is polynomial in the size of the set-SLPs.
\end{exa}


\begin{lem}
\label{lemma:fast-pair}
    Given a state $p$, a set-SLP of length $\OO(n^2)$ defining a set $\{w_q \mid q \in \Mer(p)\}$, where $w_q$ is a word with $p \in p \cdot w_q$ and $p \in q \cdot w_q$, can be computed in $\OO(mn^4)$ time and $\OO(n^2)$ space.
\end{lem}
\begin{proof}[Proof sketch]
Call vertices $(p, q)$ with $q \in \Mer(p)$ \emph{merging}. The idea is to construct, by a digraph search of $G^{(2)}$, a directed tree $T$ rooted in $(p, p)$  and containing a path from each merging vertex to the root, and then use the joint subpaths of these paths in the tree to obtain a short set-SLP describing these paths. See \autoref{fig:square} for an example.
\end{proof}
\begin{proof}[Proof of \autoref{lemma:fast-pair}]
Call vertices $(p, q)$ with $q \in \Mer(p)$ \emph{merging}. The idea is to construct a directed tree rooted in $(p, p)$ and containing paths from each merging vertex to the root, and then use the joint subpaths of these paths in the forest to obtain a short set-SLP description of these paths.

First, by performing a backwards digraph search in $G^{(2)}$ starting from $(p, p)$, find a subgraph $T$ of $G^{(2)}$ with the following properties (see \autoref{fig:square} for an illustration):
\begin{itemize}
    \item $T$ is a directed tree directed towards its root $(p, p)$;

    \item if $q \in \Mer(q)$, then $T$ contains $(p, q)$ or $(q, p)$;

    \item every leaf of $T$ is a merging vertex, and these are the only leaves in $T$.
\end{itemize}

% Fix indentation
\noindent
Call a directed path $\rho$ in $T$ a \emph{maximal branch} if all its edges belong to $T$, every vertex of~$\rho$ except the last one has outdegree exactly one in $T$, and merging vertices only occur as the first or the last vertex in $\rho$. Intuitively, maximal branches are paths in $T$ between each pair of consecutive branching, leaf, root or merging vertices of~$T$. See \autoref{fig:square} for an example.
Clearly, computing $T$ and then $\rho_1, \ldots, \rho_k$ can be done in time linear in $|E^{(2)}|$.

Let $w_1, \ldots, w_k$ be the words labeling the paths $\rho_1, \ldots, \rho_k$. 
By going backwards from the root of $T$, we can construct a set-SLP for the required words by expressing them as concatenations of $w_1, \ldots, w_k$. This can easily be done in time linear in $|E^{(2)}|$.

It remains to estimate the length of the obtained set-SLP. Since $T$ is a tree, the total length of all words $w_1, \ldots, w_k$ is at most $|V^{(2)}| = n^2$. The number of maximal branches $k$ is~$\OO(n)$. Hence, the length of the obtained set-SLP is~$\OO(n^2)$.
\end{proof}

    
\begin{figure}[h]
\centering
% Center the two drawings as a group
\begin{tabular}{@{}c@{\hspace{1cm}}c@{}}
\begin{tikzpicture} [node distance = 2cm]
%,
\tikzset{every state/.style={inner sep=1pt,minimum size=1.5em}}

\node [state] at (0, 0.5) (1) {$1$};
\node [state] at (1, 1.5) (2) {$2$};
\node [state] at (2.5, 1.5) (3) {$3$};
\node [state] at (3.5, 0.5) (4) {$4$};
\node [state] at (3.5, -1) (5) {$5$};
\node [state] at (2.5, -2) (6) {$6$};
\node [state] at (1, -2) (7) {$7$};
\node [state] at (0, -1) (8) {$8$};




\path [-stealth, thick]
(1) edge [bend left=10] node[left] {$a$} (2)
(2) edge [bend left=10] node[above] {$a, b$} (3)
(3) edge [bend left=10] node[right] {$a, b$} (4)
(4) edge [bend left=10] node[right] {$a, b$} (5)
(5) edge [bend left=10] node[right] {$a, b$} (6)
(6) edge [bend left=10] node[below] {$b$} (7)
(7) edge [bend left=10] node[left] {$a, b$} (8)
(8) edge [bend left=10] node[left] {$a, b$} (1)

(1) edge [] node[above] {$b$} (4)
(6) edge [] node[right] {$a$} (3)


;
\end{tikzpicture}
&
\begin{tikzpicture} [node distance = 2cm]
%,
\tikzset{every state/.style={inner sep=1pt,minimum size=1.5em}}

\node [state] at (-1, 0) (4) {$4, 4$};


\node [state] at (0, 0) (13) {$1, 3$};

\node [state] at (-1, -2.5) (31) {$3, 1$};
\node [state] at (1, 1) (24) {$2, 4$};
\node [state] at (2, 1) (35) {$3, 5$};
\node [state] at (3, 0) (46) {$4, 6$};
\node [state,accepting] at (3, -1) (57) {$5, 7$};
\node [state] at (2, -2) (68) {$6, 8$};
\node [state,accepting] at (1, -2) (17) {$7, 1$};
\node [state] at (0, -1) (28) {$8, 2$};

\node [state] at (4.5, 0) (15) {$1, 5$};
\node [state] at (5.5, 0) (26) {$2, 6$};
\node [state,accepting] at (5.5, -1) (37) {$7, 3$};
\node [state] at (4.5, -1) (48) {$8, 4$};

\node [state] at (5.5, 1) (3) {$3, 3$};


\path [-stealth, thick]
(13) edge [bend left=10] node[above] {} (24)
(24) edge [bend left=10] node[above] {} (35)
(35) edge [bend left=10] node[above] {} (46)

(46) edge [dotted, bend left=10] node[right] {$a$} (57)
(57) edge [bend left=10, dotted] node[below] {$a$} (68)
(68) edge [bend left=10] node[above] {} (17)
(17) edge [bend left=10, dotted] node[below] {$a$} (28)
(28) edge [bend left=10, dotted] node[left] {$a$} (13)


(13) edge [dotted] node[above] {$b$} (4)
(68) edge [] node[above] {} (13)

(68) edge [dotted, bend left=20] node[below] {$a$} (31)
(31) edge [dotted] node[left] {$b$} (4)

(15) edge [] node[above] {} (26)
(26) edge [] node[above] {} (37)
(37) edge [dotted] node[below] {$a$} (48)
(48) edge [dotted] node[left] {$a$} (15)

(15) edge [dotted] node[above] {$b$} (46)

(26) edge [] node[above] {} (3)

;
\end{tikzpicture}
\end{tabular}
\caption{An example of $\Aa$ (left), and a part of the underlying digraph $G^{(2)}$ of its square automaton (right). Merging vertices are doubly circled. The edges of $T$ for $p = 7$ are represented by dotted edges, and these edges are labelled with one of the letters labelling the corresponding transition in $\Aa^{(2)}$. Furthermore, we have $\rho_1 = (7, 1) \to (8, 2) \to (1, 3) \to (4, 4)$, $\rho_2 = (5, 7) \to (6, 8) \to (3, 1) \to (4, 4)$, $\rho_3 = (7, 3) \to (8, 4) \to (1, 5) \to (4, 6) \to (5, 7)$. A set-SLP encoding the labels of these path is presented in \autoref{ex:set-slp}, with~$u_i$ labelling the path $\rho_i$, $i \in \{1, 2, 3\}$.
}\label{fig:square}
\end{figure}

% Fix indentation
\noindent
To construct a maximal column more efficiently, we will use \cite[Corollary 11]{Chistikov2026Arxiv}.

% propC: cite-only heading (avoid double parentheses)
\begin{propC}[\cite{Chistikov2026Arxiv}]\label{prop:mat-mult}
    Given $M_1, M_2 \in \{0, 1\}^{n \times n}$ such that $M_1 M_2 \in \{0, 1\}^{n \times n}$, one can compute $M_1 M_2$ in $\OO(n^2)$ time.
\end{propC}

% Fix indentation
\noindent
Since in the remainder of this section we are dealing only with matrices from a fixed zero-one matrix monoid, \autoref{prop:mat-mult} means that we can always perform matrix multiplication in quadratic time.

\begin{lem}\label{lemma:time-max-col}
    For a given state $p$ of $\Aa$, an SLP of length $\OO(n^2)$ encoding a word $w$ such that $[w \cdot p]$ is a maximal column can be computed in $\OO(mn^4)$ time and $\OO(n^3)$ space.
\end{lem}
\begin{proof}[Proof sketch]
Compute the matrices of the words encoded by the set-SLP from \autoref{lemma:fast-pair}. 
We rely on the property, already used in some form in \cite{Carpi1988}, that if for all $q \in \Mer(p)$, $q \ne p$, the vector $[w \cdot q]$ is zero, then $[w \cdot p]$ is a maximal column. To construct a word with this property, we iteratively concatenate the words $w_q$ depending on nonzero columns in the matrix in each iteration. The number of iterations is bounded by $n$.
\end{proof}

\begin{proof}[Proof of \autoref{lemma:time-max-col}]
By \autoref{lemma:fast-pair}, 
we can compute in $\OO(mn^4)$ time a set-SLP of length $\OO(n^2)$ defining a set $\{w_q \mid q \in \Mer(p)\}$, where $w_q$ is a word with $p \in p \cdot w_q$ and $p \in q \cdot w_q$. For each~$q \in \Mer(p)$, we then compute $M(w_q)$, which in total requires $\OO(n^2)$ matrix multiplications and $\OO(n^3)$ space.

Perform now the following algorithm that iteratively constructs SLPs for words $w_i$ and~$M(w_i)$ for~$i \ge 0$. Let $w_0$ be the empty word. Assume that an SLP for $w_{i - 1}$ and $M(w_{i - 1})$ are already constructed. If there is a state $q \in \Mer(p)$, $q \ne p$, such that $[w_{i - 1} \cdot q] = M(w_{i - 1})[q]$ is not the zero vector, take $w_i = w_{i - 1} w_q$ and compute $M(w_i)$, otherwise stop and output~$w_{i - 1}$. Clearly, each step of the algorithm only requires a constant number of matrix multiplications.

The algorithm clearly terminates in at most $|\Mer(p)|$ steps, since each iteration decrements by at least one the number of states $q \in \Mer(p)$ such that $[w_i \cdot q]$ is not the zero vector. It is also easy to see that the constructed SLP for $w$ has length $\OO(n^2)$.

Assume now that $[w \cdot p]$ is not a maximal column. Let $C \subseteq Q$ be such that $[C] \ge [w \cdot p]$ is a maximal column. By definition of $\Mer(p)$, $C \subseteq \Mer(p)$, hence for every $q \in C$, $q \ne p$, the vector $[w \cdot q]$ is zero. Since $[C] \ge [w \cdot p]$, $p \in C$. Thus we get that $M(w)[C] = \sum_{q \in C} [w \cdot q] = [w \cdot p]$.
By \autoref{lem:max-column-stable}, $M(w)[C]$ is a maximal column, hence $[w \cdot q]$ must also be a maximal column. 
\end{proof}



\subsection{Finding a matrix of minimum rank}\label{subsec-finding-mat}

We now use the results of the previous section to construct a matrix of minimum rank. The key idea is as follows: since the set of maximal columns is stable under left multiplication by matrices from the monoid, we can iteratively make each column of the matrix maximal or zero by, intuitively, applying the same word (together with a short ``reachability'' word) to a state in each column. This simple observation significantly decreases the time complexity of our algorithm compared to the naive implementation of the algorithm from \cite{Carpi1988} constructing a word of minimum rank. Indeed, in the approach of \cite{Carpi1988}, the word is constructed letter by letter, and requires to know the result of applying the last letter at every step. Since the constructed word has length $\OO(rn^3) = \OO(n^4)$, where $n$ is the number of states and $r$ is the rank, this results in $\OO(n^{4+\omega})$ time complexity. By efficiently constructing only one maximal column (as described in the previous section) and reusing it for the whole set of states (as described in this section), we decrease the time complexity to $\OO(n^{2+\omega})$.

\begin{prop}\label{prop:time-minrank-mat}
    An SLP of length $\OO(n^2)$ encoding a word $w$ of minimum rank can be computed in $\OO(mn^4)$ time and $\OO(n^3)$ space.
\end{prop}
\begin{proof}[Proof sketch]
 Compute the matrix of the word $w$ encoded by the SLP from \autoref{lemma:time-max-col}. Iteratively, for each $q \in Q$, concatenate $w$ with a word of length at most $n$ mapping $p$ to a state corresponding to a nonzero element of the row in the current iteration. Denote by $w_n$ the resulting word, which has the property that all nonzero columns of $M(w_n)$ are maximal. Symmetrically compute $w'_n$ for rows. Then by \autoref{lem:all-max-column} the word $w_{n}w'_{n}w_{n}w'_{n}$ matrix has minimum rank.
\end{proof}
\begin{proof}[Proof of \autoref{prop:time-minrank-mat}]
By \autoref{lemma:time-max-col},
for a given state $p$ of $\Aa$, we can compute in $\OO(mn^4)$ time an SLP of length $\OO(n^2)$ defining a word $w$ such that $[w \cdot p]$ is a maximal column. Compute $M(w)$, which can be done in  $\OO(n^4)$ time.

Let $Q = \{q_1, \ldots, q_n\}$. Define $v_0 = \epsilon$. Clearly, $M(v_0)$ is the identity matrix. 
For each $1 \le i \le n$, perform the following algorithm. If $[v_{i - 1} \cdot q_i]$ is the zero vector, take $v_i = v_{i - 1}$ and go the the next step. Otherwise, find $q$ such that $q \in v_{i - 1} \cdot q_i$, and 
find a word $u_{p \to q}$ of length at most $n - 1$ mapping $p$ to $q$, and compute $M(u_{p \to q})$, which can be done by~$\OO(n)$ matrix multiplications. Take $w_i = w u_{p \to q} w_{i - 1}$, and compute $M(w_i)$, which requires a constant number of matrix multiplications. Go to the next step of the algorithm.

Observe that after the $i$th step of the algorithm, $[v_j \cdot q_j]$ is a maximal column. Indeed, $[v_i \cdot q_i] = [w \cdot p]$, and for all $j < i$, $[v_j \cdot q_j]$ is a result of left multiplication of a maximal column by a matrix from the monoid, which is a maximal column by \autoref{lem:max-column-stable}.
Moreover, the obtained SLP of $w_{n}$ has length $\OO(n^2)$.

It remains to construct symmetrically $w'_{n}$ for rows instead of columns. Then each column and each row of $M(w_{n}w'_{n})$ is either maximal or zero, and by \autoref{lem:all-max-column} the word $w_{n}w'_{n}w_{n}w'_{n}$ matrix has minimum rank.
\end{proof}


% Fix indentation
\noindent
Given an SLP of length $\OO(n^2)$ encoding a word $w$, we can compute the matrix of $w$ by computing the matrices of words occurring in the derivation of $w$ from bottom to top in time $\OO(n^4)$.
Thus we prove \autoref{thm:finding-matrix}. We also get the following result, which can be seen as a proof of a very weak version of the \v{C}ern\'{y} conjecture generalised from rank one words in total DFAs to minimum rank words in complete UFAs.

\begin{thm}
    For every $n$-state complete UFA, there exists an SLP of length $\OO(n^2)$ encoding a word of minimum rank.
\end{thm}

% Fix indentation
\noindent
We remark that the length of the word encoded by the constructed SLP asymptotically matches the best known upper bound  for words of minimum rank: $\OO(n^4)$ for complete UFAs~\cite{Carpi1988} and $\OO(n^3)$ for total DFAs~\cite{Klyachko1987}. In particular, one can efficiently compute words of minimum rank within these bounds.

\subsection{Total DFAs}\label{subsec-dfas}

For total DFAs, we follow the same algorithms as in the proof of \autoref{thm:finding-matrix}, but exploit the fact that elementary matrix operations can be performed more efficiently. Namely, if $\Aa$ is a total DFA, then each word defines a transformation on $Q$. By storing matrices of words as transformations, we get that matrix multiplication can be performed in $\OO(n)$ time, and each matrix requires $\OO(n)$ space. Moreover, we have $|E^{(2)}| = mn^2$. By taking these improvements into account, we get the proof of \autoref{thm:main-dfas}.


